% 长度外推的两个维度；位置与块数仅作教学示意。
\begin{tikzpicture}[x=1mm,y=1mm,font=\kaishu\footnotesize]
 \node[anchor=west,text=BookInk] at (0,49){输入更长：在更多信息中定位和组合证据};
 \foreach \x in {2,13,24,35,46,57,68,79,90}{
  \draw[BookRule,fill=BookPaper] (\x,36) rectangle +(9,8);
 }
 \fill[BookTeal!30] (13,36) rectangle +(9,8);
 \fill[BookTeal!30] (68,36) rectangle +(9,8);
 \node[font=\kaishu\scriptsize] at (17.5,40){证据};
 \node[font=\kaishu\scriptsize] at (72.5,40){证据};
 \draw[book arrow,BookTeal] (18,35)--(18,30)--(72,30)--(72,35);
 \node[below,text=BookMuted,font=\kaishu\scriptsize] at (45,30){位置、间距和干扰项都可能改变};
 \node[anchor=west,text=BookInk] at (0,18){步骤更多：把局部运算连续执行更久};
 \foreach \x/\t in {2/1,24/2,46/3,83/K}{
  \draw[BookBlue,fill=BookBlue!7] (\x,4) rectangle +(14,8);
  \node at (\x+7,8){步骤$\t$};
 }
 \draw[book arrow] (16,8)--(23,8);
 \draw[book arrow] (38,8)--(45,8);
 \draw[book arrow] (60,8)--(66,8);
 \node at (73,8){$\cdots$};
 \draw[book arrow] (79,8)--(82,8);
 \node[below,text=BookMuted,font=\kaishu\scriptsize] at (50,1){中间状态改变，单步错误可能累积或放大};
\end{tikzpicture}
